\begin{theorem}[Finite Lebesgue-number Phase Real source lock]
\label{thm:finite-lebesgue-number-phase-real-source-lock}
Every accepted public radius read of $\FiniteLebesgueNumberUp$ that is visible to a compact or continuous consumer factors through the displayed Phase Real and Completion route
$$
\DyadicCoverUp,\ \RationalBallUp,\ \DyadicMeshUp,\ \StreamNameUp,\ \RegSeqRatUp,\ \RealUp .
$$
Thus the consumer-facing radius is not a host scalar and not an open-cover compactness witness: it is the finite dyadic-radius packet carried by the $\BHist$, $\hsame$, $\Cont$, $\Pkg$, and $\NameCert$ rows of the finite Lebesgue-number carrier before any $\CompactMetricUp$, $\TotallyBoundedUp$, $\CompactNetModulusSelectorUp$, $\UniformModulusUp$, or $\ContinuousMapUp$ endpoint can inspect it.
\end{theorem}
\leanchecked{BEDC.Derived.FiniteLebesgueNumberUp.FiniteLebesgueNumberPhaseRealSourceLock}
\begin{proof}
Start with the root obligations. \autoref{thm:finite-lebesgue-number-namecert-obligations} fixes the dyadic cover, finite-window, rational-ball radius, mesh, transport, continuation, provenance, and local naming rows as the only public rows of the $\FiniteLebesgueNumberUp$ packet.

Next apply the L10 source lock. \autoref{thm:finite-lebesgue-number-l10-source-lock} forces the radius route through $\StreamNameUp$, $\RegSeqRatUp$, $\RealUp$, $\DyadicMeshUp$, and $\DyadicCoverUp$ before the compact side can read the radius.

Use the selected-tail bridge ledger and bridge packet exactness for the Phase Real and Completion handoff. \autoref{def:finite-lebesgue-number-selected-tail-bridge-ledger} records the finite selected-tail source rows, while \autoref{thm:finite-lebesgue-number-bridge-packet-exactness} leaves no bridge coordinate outside the displayed dyadic, stream, regular-rational, and real rows.

Finally use compact-continuous unblock. \autoref{thm:finite-lebesgue-number-compact-continuous-triad} sends compact, compact-net, continuous-map, and uniform-modulus consumers through the same finite dyadic-radius packet. Hence every accepted public read factors through the stated Phase Real and Completion route.
\end{proof}

\begin{theorem}[Finite Lebesgue-number selected-tail modulus exhaustion]
\label{thm:finite-lebesgue-number-selected-tail-modulus-exhaustion}
In an accepted $\FiniteLebesgueNumberUp$ route, every selected-tail witness and every finite-window witness used to expose a radius to a modulus consumer is exhausted by the displayed carrier rows
$$
D,W,\lambda,M,A_T,W_T,K_T,S_T,R_T,H,C,P,N .
$$
No hidden modulus selector, quotient stream, countable choice coordinate, host real equality, or detached $\RealUp$ representative can be inserted between the selected-tail source rows and the finite dyadic-radius packet.
\end{theorem}
\leanchecked{BEDC.Derived.FiniteLebesgueNumberUp.FiniteLebesgueNumberSelectedTailModulusExhaustion}
\begin{proof}
Read the finite-window side first. \autoref{thm:finite-lebesgue-number-namecert-obligations} and \autoref{thm:finite-lebesgue-number-l10-source-lock} make every window and radius read pass through the dyadic cover, rational-ball radius, mesh, stream-name, regular-rational, and real rows with the displayed structural support.

Now read the selected-tail side. \autoref{def:finite-lebesgue-number-selected-tail-bridge-ledger} lists exactly the selected finite-tail admission, stream-window, regular-rational readback, dyadic tolerance, and Real-seal rows joined to the finite Lebesgue-number packet.

Apply bridge packet exactness. \autoref{thm:finite-lebesgue-number-bridge-packet-exactness} states that the bridge packet has no coordinate beyond the displayed tuple. The excluded coordinates include quotient equality, host real equality, choice, detached real completeness, and external modulus data.

The modulus consumer therefore sees only the selected-tail and finite-window rows already named in the carrier. Since there is no further selector or stream coordinate in the packet, the listed rows exhaust the modulus-facing witness data.
\end{proof}

\begin{theorem}[Finite Lebesgue-number compact-continuous nonchoice exit]
\label{thm:finite-lebesgue-number-compact-continuous-nonchoice-exit}
Any $\CompactMetricUp$, $\TotallyBoundedUp$, $\CompactNetModulusSelectorUp$, $\UniformModulusUp$, or $\ContinuousMapUp$ consumer admitted downstream of $\FiniteLebesgueNumberUp$ receives exactly the finite dyadic-radius packet
$$
\DyadicCoverUp,\ \RationalBallUp,\ \DyadicMeshUp,\ \StreamNameUp,\ \RegSeqRatUp,\ \RealUp,\ \BHist,\ \hsame,\ \Cont,\ \Pkg,\ \NameCert .
$$
The exit to these consumers does not supply open-cover compactness, an arbitrary subcover principle, a host minimum, countable choice, quotient stream equality, or an external modulus selector.
\end{theorem}
\leanchecked{BEDC.Derived.FiniteLebesgueNumberUp.FiniteLebesgueNumberCompactContinuousNonchoiceExit}
\begin{proof}
Begin with source accountability. \autoref{thm:finite-lebesgue-number-l10-source-lock} and \autoref{thm:finite-lebesgue-number-phase-real-source-lock} identify the only source of the consumer-facing radius as the finite dyadic, stream, regular-rational, and real route.

Next use the selected-tail exhaustion theorem above. It prevents selected-tail or finite-window witnesses from contributing a hidden modulus selector, quotient stream, countable choice coordinate, host real equality, or detached real representative.

Finally apply the downstream consumer package. \autoref{thm:finite-lebesgue-number-compact-continuous-triad} sends compact, compact-net, continuous-map, and uniform-modulus endpoints through the same finite packet, and \autoref{thm:finite-lebesgue-number-bridge-packet-exactness} removes all bridge coordinates outside that packet.

Therefore the consumer exit is nonchoice: the consumers receive the finite dyadic-radius packet and no open-cover compactness or arbitrary subcover principle.
\end{proof}

\begin{theorem}[Finite Lebesgue-number real-phase coverage export]
\label{thm:finite-lebesgue-number-real-phase-coverage-export}
Every accepted $\FiniteLebesgueNumberUp$ public read used by a $\CompactMetricUp$, $\CompactNetModulusSelectorUp$, $\ContinuousMapUp$, or $\UniformModulusUp$ endpoint factors through exactly the L10 Real-phase route
$$
\DyadicRatCoreUp,\ \StreamNameUp,\ \RegSeqRatUp,\ \RealUp
$$
together with the finite dyadic cover, rational-ball radius, dyadic mesh, $\BHist$, $\hsame$, $\Cont$, $\Pkg$, and $\NameCert$ rows of the finite Lebesgue-number carrier. The exported handoff supplies no host real equality, quotient stream equality, countable-choice coordinate, selected limit, open-cover compactness principle, or ambient-completeness primitive.
\end{theorem}
\leanchecked{BEDC.Derived.FiniteLebesgueNumberUp.FiniteLebesgueNumberRealPhaseCoverageExport}
\begin{proof}
Start with the Phase Real source lock. \autoref{thm:finite-lebesgue-number-phase-real-source-lock} routes every compact- or continuous-visible radius read through the dyadic cover, rational-ball radius, dyadic mesh, stream-name, regular-rational, and Real rows, with $\BHist$, $\hsame$, $\Cont$, $\Pkg$, and $\NameCert$ carrying the public BEDC packet.

Next use selected-tail exhaustion. \autoref{thm:finite-lebesgue-number-selected-tail-modulus-exhaustion} rules out any hidden modulus selector, quotient stream, countable choice coordinate, host real equality, or detached $\RealUp$ representative between the selected-tail source rows and the finite dyadic-radius packet.

Finally apply the compact-continuous nonchoice exit. \autoref{thm:finite-lebesgue-number-compact-continuous-nonchoice-exit} says that the compact, compact-net, continuous-map, and uniform-modulus endpoints receive exactly that finite packet and no open-cover compactness, arbitrary subcover principle, host minimum, quotient stream equality, or external modulus selector.

Thus the L10 handoff exported to the downstream endpoints is the finite Real-phase coverage route through $\DyadicRatCoreUp$, $\StreamNameUp$, $\RegSeqRatUp$, and $\RealUp$, together with the named finite carrier rows and no completed-real or choice-bearing primitive.
\end{proof}
